\clearpage

\section {Introduction }
\label{sec:intro}

In this technical note, 
we report the results of a 3-flavour $\nu_{\mu}$-disappearance analysis,
performed by the Valencia-Lancaster-Oxford-Rutherford (VaLOR) group,        
on the combined Run1+2+3 (3.010$\times$10$^{20}$ POT) dataset.
This dataset was collected in three run periods:
in January 2010 - June 2010 (Run1, 0.323$\times$10$^{20}$ POT),
in November 2010 - March 2011 (Run2, 1.108$\times$10$^{20}$ POT) and,
following the recovery work from the damage of the 2011 Tohoku earthquake, 
in March 2012 - June 2012 (Run3, 1.579$\times$10$^{20}$ POT).

This analysis is an update of the work presented in
Refs. \cite{OurRun1DisappAnalysis}, \cite{OurRun12DisappAnalysis} and \cite{OurRun12DisappAnalysis3Flavour}
with improved methods and an enlarged dataset. It benefits from an updated flux tuning based on
recent NA61 data \cite{FluxTuning}\cite{NA61T2K}, as well as from T2K fits of flux and cross-section parameters to
ND280 measurements of muon momentum-angle distributions of $\nu_\mu$ CCQE-like and CCnonQE-like events whose 
vertices are in the FGD \cite{BANFF}.
Details of our updated method for predicting the reconstructed energy spectrum 
of single $\mu$-like ring events in SuperK are presented in Sec.\ \ref{sec:spectrum_prediction}.

The oscillation analysis is performed using neutrino oscillation probabilities computed in a 3-flavour
framework including constant-density matter effects. 
Details of our oscillation probability calculations, estimates of the numerical accuracy of our
calculations and comparisons with calculations used by the alternative T2K oscillation analyses, 
are given in Appendices \ref{appendix:valor_3f_osc_prob_code} and \ref{appendix:comparison_with_prob3}.

We consider a total of 48 efficiency, energy scale, neutrino flux, neutrino cross-section, 
initial-state nuclear environment and final-state rescattering systematics. 
These systematic parameters and their effects 
on the extrapolated single $\mu$-like ring event reconstructed energy spectrum are described in Sec.\ \ref{sec:systematics}.
Further details of them are also given in Appendices \ref{appendix:nsk_systematics} and \ref{appendix:spectrum_systematics}.

As in our previous analyses, the oscillation fit was made using the likelihood-ratio method \cite{StatAnalysisCowan}.
The oscillation fitting method is presented in detail in Sec.\ \ref{sec:oscfit_method}, while the chosen
goodness-of-fit test, suitable for the Run 1+2+3 statistics, is described in Sec.\ \ref{sec:gof_test}.
The fitter performance was tested in detail using toy MC experiments 
and the results of these tests are presented in Sec.\ \ref{sec:toy_mc_fitter_performance} 
and in Appendix \ref{sec:appendix:fitter_performance}.
The effect of systematic parameters on the oscillation fit was studied in detail. 
The results of this study are summarised in Sec.\ \ref{sec:fit_syst_analysis} and are shown, 
in more detail, in Appendix \ref{appendix:systematics_bestfit_oscillations}.

We considered two versions of the fit: the first version allows all 48 systematic parameters to float 
and is denoted the `2+N' fit, while the second version allows only the 8 most dominant systematics to float 
(the `2+8' fit). These 8 dominant systematic parameters are those having an effect on the oscillation
parameters greater than 25\% of the expected statistical error, for at least one set of true oscillation
parameters in the parameter range of interest. Results are presented both for `2+N' and `2+8' fits.
Further details of our fit labeling scheme are given in Sec.\ \ref{sec:oscfit_method_label}.

The confidence regions were constructed, as in our previous analyses, 
using the Feldman-Cousins method \cite{FeldmanCousins} both for `2+N' and `2+8' fits. 
Systematic errors were included in the oscillation contours using the 
Cousins-Highland method (Bayesian integration) \cite{CousinsHighland}.
Contours constructed with the constant-$\Delta\chi^{2}$ method are also shown for both fits for comparison.
In Sec.\ \ref{sec:contour_construction} we summarise both of these methods
and clearly explain the exact methodology used for incorporating systematic errors.

Sensitivity studies are shown in Appendix\ \ref{appendix:sensitivity_study}, and were performed using the complete VaLOR analysis 
with an integrated exposure of 3.010$\times$10$^{20}$ POT and for a grid of oscillation points around the 2012 MINOS \cite{MINOS2012} 90\% CL region.

The results of the measurement of $\sin^{2}2\theta_{23}$ / $\sin^{2}2\theta_{23}$ and $|\Delta m^{2}_{32}|$ 
using the Run1+2+3 dataset are given in Sec.\ \ref{sec:results_bestfit} of this note for both `2+N' and `2+8' fits,
while the goodness-of-fit tests are presented in Sec.\ \ref{sec:results_gof}.
The best-fit spectra for all `2+N' and `2+8' fits to the Run 1+2, Run 3 and Run 1+2+3 datasets are shown in
Sec.\ \ref{sec:results_bestfit_spectra}.
The Feldman-Cousins and constant-$\Delta\chi^{2}$ confidence regions are shown, for both fits, 
in Sec.\ \ref{sec:results_contours}. Comparisons of our results obtained with different fits and with
different methods for constructing confidence regions are presented in Sec.\ \ref{sec:results_contour_comparisons}.
In that section, we also show comparisons of our T2K Run 1+2+3 results with MINOS and SuperK results, 
as well as with the expected T2K sensitivity for an integrated exposure of 3.010$\times$10$^{20}$ POT.
We present the VaLOR group prespective on the choice 
of primary result for the Run1+2+3 disappearance analysis in Sec.\ \ref{sec:primary_result_selection}.
Finally, a summary is given in Sec.\ \ref{sec:summary}.  

